\section{Introduction}\label{sec:introduction}

Meeting delivery commitments requires manufacturers to coordinate machine
assignment, sequencing, and completion timing. Finishing a job too early may
create inventory and holding costs, while finishing it too late may cause service
failures or contractual penalties. Earliness--tardiness scheduling formalizes
this trade-off \citep{baker1990review,kramer2019unified}; a due window further
recognizes that completion is acceptable over an interval rather than at one
prescribed instant \citep{wan2002tabu,rosa2017algorithms}. The resulting decision
is not simply to process jobs as early as possible. A manufacturer must decide
which jobs share a machine, in what order they are processed, and where deliberate
idle time should be inserted to align their completion times with different
delivery windows.

Internal production is not the only way to meet these commitments. Subcontracting
provides a capacity option when the internal system is congested or when processing
a particular job internally is uneconomical \citep{vanmieghem1999subcontracting}.
In a scheduling setting, however, outsourcing and internal production cannot be
optimized independently. Outsourcing one job removes its processing and setup
requirements from an internal machine, changes the feasible positions of the
remaining jobs, and may reduce both earliness and tardiness costs. Existing studies
therefore jointly decide which jobs are outsourced and how the retained jobs are
scheduled in single-machine, flow-shop, job-shop, or coordinated
production--distribution systems
\citep{lee2008outsourcing,safarzadeh2019jobshop,kim2022flowshop,zhong2022coordinated}.

A further coupling arises when the external provider prices the outsourced
portfolio rather than each job separately. Quantity discounts or spend-based
tariffs make the payment depend on the aggregate baseline cost of all outsourced
jobs \citep{lu2021discount}. The marginal outsourcing cost of one job then depends
on which other jobs are outsourced with it. Similar portfolio pricing has also
motivated alternative exact decompositions in private-fleet/common-carrier
planning \citep{dabia2019rich}. In the present scheduling setting, this price
coupling interacts with due-window congestion: outsourcing a group of jobs can
simultaneously move the total payment to a different tariff range and reshape the
internal machine schedules.

The relevant scheduling literature provides these ingredients only separately.
Due-window research models general early and late completion costs, including
sequence-dependent setups, but normally requires all jobs to be processed
internally \citep{wan2002tabu,rosa2017algorithms}. Parallel-machine exact methods
provide strong relaxations and column-generation frameworks for machine schedules
\citep{kedad2008lowerbounds,sen2015preemptive,bulhoes2020exact}. Scheduling with
outsourcing or rejection decides which jobs leave the internal system, but most
models use additive job-level charges, outsourcing lead times, or objectives that
do not optimize job-dependent due-window costs
\citep{slotnick2011review,lee2008outsourcing,kim2022flowshop,zhong2022coordinated}.
To the best of our knowledge, no exact approach jointly addresses identical
parallel machines, job-dependent due windows, sequence-dependent internal setups,
and a nonseparable aggregate outsourcing tariff.

This combination also exposes an algorithmic limitation. In a Dantzig--Wolfe
decomposition, an internal column is a complete single-machine schedule. For
integer-time data, its pricing problem can be solved on a time-indexed acyclic
graph \citep{bulhoes2020exact}, but the numbers of states and arcs grow with the
numerical time horizon.
Moreover, if elementarity is relaxed in this graph, negative non-elementary
pseudo-schedules may be found even though only elementary schedules can enter the
master problem. A continuous-time alternative avoids the explicit horizon
expansion, but it cannot value a partial sequence by one scalar: the best reduced
cost depends on the completion time at its boundary. This calls for a functional
labeling method that controls repeated visits while preserving an exact pricing
certificate.

We address these gaps through the following contributions.
\begin{enumerate}[label=(\roman*)]
    \item We introduce an integrated parallel-machine due-window scheduling
    problem with aggregate outsourcing. The model couples job sourcing with
    internal assignment, sequencing, timing, and sequence-dependent setups.

    \item We derive two set-covering formulations that place the outsourcing
    decision on different sides of the decomposition. Formulation \textup{(SP1)}
    represents outsourced job sets by columns, whereas formulation \textup{(SP2)}
    keeps job-level outsourcing and tariff variables in the master problem.
    This distinction accommodates different outsourcing cost structures and
    exposes the trade-off between a second pricing family and a richer master.

    \item We develop an exact continuous-time internal pricing method based on
    bidirectional piecewise-linear labeling and ng-route decremental state-space
    refinement. Its source-aware dominance structure maintains collective
    functional envelopes, and the method does not explicitly discretize the time
    horizon. An exact time-indexed method is retained as an alternative for
    integer-time instances.

    \item We embed both formulations and both internal pricing approaches in a
    unified branch-price-and-cut framework with formulation-specific outsourcing
    pricing, branching, certified node bounds, and Phase-I feasibility repair.
\end{enumerate}

The remainder of the paper is organized as follows.
Section~\ref{sec:literature} reviews due-window scheduling and scheduling with
outsourcing. Section~\ref{sec:problem} defines the integrated decision problem and
its cost structure. Section~\ref{sec:formulations} presents formulations
\textup{(SP1)} and \textup{(SP2)}. Section~\ref{sec:solution} describes the exact
solution framework and the internal and outsourcing pricing algorithms.
Section~\ref{sec:experiments} reports the computational study, and
Section~\ref{sec:conclusions} concludes.
